% §5 Problem 2: coordinates of the looping case, Lagrangian against particle simulation
% Data: Data/S5P2-verification.csv, written by python Simulations/study.py S5P2
\documentclass{SolutionFigure}

\begin{document}
\begin{tikzpicture}
    \begin{groupplot}[
        subplot grid, % two 6 x 4 cm panels
        group style={ylabels at=all, horizontal sep=20mm}, % different units
        xlabel={$t$ (s)}, xmin=0, xmax=5,
        table/col sep=comma,
    ]
        % Lagrangian curves, particle markers
        \nextgroupplot[title={(a) Support}, ylabel={$x$ (m)}, ymin=-0.5, ymax=1.5,
            legend pos=north west, legend columns=2,
            legend style={/tikz/every even column/.append style={column sep=2mm}}]
            \addplot table[x=t, y=x] {Figures/Data/S5P2-verification.csv};
            \addlegendentry{Lagrangian}
            \addplot[plot points, mark size=0.45mm, each nth point=4]
                table[x=t, y=x-engine] {Figures/Data/S5P2-verification.csv};
            \addlegendentry{particles}

        \nextgroupplot[title={(b) Angle}, ylabel={$\phi$ (rad)}, ymin=0, ymax=25]
            \addplot table[x=t, y=phi] {Figures/Data/S5P2-verification.csv};
            \addplot[plot points, mark size=0.45mm, each nth point=4]
                table[x=t, y=phi-engine] {Figures/Data/S5P2-verification.csv};
    \end{groupplot}
\end{tikzpicture}
\end{document}
